\section{Project GR00T: modelo fundacional general para robots}

\subsection{Visión y arquitectura}

El objetivo de GR00T es construir un modelo fundacional general capaz de operar diversos cuerpos robóticos:

\begin{itemize}
    \item Recibe instrucciones en lenguaje natural y entradas visuales
    \item Genera secuencias de acciones robóticas
    \item Generaliza entre distintas morfologías de robots: el mismo modelo controla robots diferentes
\end{itemize}

\subsection{Retos técnicos centrales}

\begin{warningbox}{Retos particulares de los modelos fundacionales para robots}
\begin{itemize}
    \item \textbf{Escasez de datos}: a diferencia del lenguaje, que cuenta con datos a escala de internet, los datos de operación robótica son extremadamente escasos
    \item \textbf{sim-to-real gap}: las políticas entrenadas en simulación suelen fallar en el mundo real
    \item \textbf{Seguridad}: los robots actúan directamente sobre el mundo físico, y un error puede causar daño físico
    \item \textbf{Tiempo real}: la interacción física exige una latencia de decisión del orden de milisegundos
\end{itemize}
\end{warningbox}

\subsection{Estrategias de generación de datos}

\begin{itemize}
    \item Utiliza simulación física a gran escala (Isaac Sim) para generar datos de entrenamiento
    \item Genera automáticamente tareas diversas y reward function mediante LLM (método Eureka)
    \item Aprende operaciones humanas a partir de videos de internet (aprendizaje por observación)
    \item Human demonstration + teleoperation para recopilar datos reales de operación
\end{itemize}

\subsection{Resumen de la sección}

GR00T representa la ambición de extender el paradigma de los modelos fundacionales del lenguaje al mundo físico; los retos centrales están en los datos, la transferencia sim-to-real y la seguridad.
